\section{Related Work}
\label{sec:related_work}

\paragraph{Code World Models and Execution Simulation.}
Code World Models~\cite{meta2025cwm} mid-train a 32B LLM on execution traces to predict output strings and establish the generation-based paradigm for execution-free verification.
Self-Execution~\cite{maimon2026selfexecution} extends this with step-by-step trace prediction.
CodeMind~\cite{liu2024codemind} evaluates code reasoning along multiple dimensions including execution prediction.
SolidCoder~\cite{lee2026solidcoder} bridges mental simulation and actual execution via concrete feedback.
Abdollahi et al.~\cite{abdollahi2025demystifying} identify systematic failure modes in LLM trace prediction, including position-dependent errors.
All share a generation formulation whose cost scales with output length. We provide the first controlled comparison of this formulation against classification and trajectory-level alternatives at matched model capacity.

\paragraph{Execution-free Verification.}
SWE-RM~\cite{shum2025swerm} trains a trajectory-level reward model for patch ranking without execution.
Patch Reasoner~\cite{xu2025patchreasoner} uses chain-of-thought reasoning for patch verification.
LEVER~\cite{ni2023lever} trains a classifier to re-rank code candidates by predicting execution correctness; CodeT~\cite{chen2023codet} uses dual execution agreement as a verification signal.
ICE-Score~\cite{zhuo2024icescore}, CodeScore~\cite{dong2024codescore}, and CodeBERTScore~\cite{zhou2023codebertscore} evaluate code quality without execution via learned metrics, but target generation quality assessment rather than per-test pass/fail verification.
Snell et al.~\cite{snell2024scaling} show that test-time scaling via verification outperforms parameter scaling; S*~\cite{li2025sstar} and VG-Search~\cite{chen2025vgsearch} find that finer-grained verification improves selection efficiency.
Self-consistency~\cite{wang2023selfconsistency} uses majority voting over sampled solutions as an implicit verifier.
These works treat verification granularity and formulation as design choices but do not compare them at controlled capacity. We isolate the formulation variable (classification vs.\ generation vs.\ trajectory-level) while holding model, data, and compute constant.

\paragraph{Process vs.\ Outcome Supervision.}
Lightman et al.~\cite{lightman2023letsverify} demonstrate that process-based reward models (verifying each reasoning step) outperform outcome-based models (verifying only the final answer) for mathematical reasoning.
Uesato et al.~\cite{uesato2022solving} show similar advantages for process feedback in math word problems.
In code, DreamPRM-Code~\cite{dreamprm2025code} applies function-level process rewards, and PRL-Coder~\cite{ye2025prlcoder} integrates process supervision into RL-based code generation.
Our per-test classification is structurally analogous to process supervision: each test verdict corresponds to an intermediate verification step, while trajectory-level scoring corresponds to outcome supervision. Our finding that per-test and trajectory-level achieve equivalent accuracy (TOST $p = 0.007$) with per-test providing additional diagnostic value parallels the process supervision result that step-level verification enables error localization at no accuracy cost.

\paragraph{Verification in Repair Pipelines.}
CodeRL~\cite{le2022coderl} trains a critic model to provide execution feedback for RL-based code generation.
AlphaCode~\cite{li2022alphacode} and AlphaCode 2~\cite{alphacode2023team} use massive sampling with test-based filtering for competition programming.
Reflexion~\cite{shinn2023reflexion} uses verbal self-reflection on execution results to iteratively improve agent performance.
SWE-bench~\cite{jimenez2024swebench} and SWE-bench Verified~\cite{chowdhury2024swebenchverified} provide standard benchmarks; SWE-agent~\cite{yang2024sweagent} and Agentless~\cite{xia2024agentless} demonstrate agent-based and pipeline-based repair respectively.
Yang et al.~\cite{yang2025aprsurvey} survey LLM-based repair, identifying verification as a persistent bottleneck with cost scaling linearly in candidates.
These systems demonstrate that verification quality is a primary bottleneck for iterative repair. Our formulation comparison is complementary: the efficiency and diagnostic granularity of the verifier directly affects the viability of verification-in-the-loop approaches at scale.

\paragraph{Predictive Mutation Testing.}
Predicting test outcomes for code mutations without execution has been studied in software engineering~\cite{tufano2019predictive,kim2022predictive}, with recent work applying LLMs to mutation testing~\cite{wang2024llmmutation} and contextual prediction~\cite{jain2023contextual}. Our per-test classification shares this problem structure: given a code change and a test, predict pass or fail. We extend this line to LLM-based verification at repository scale, comparing it against generation and trajectory-level formulations under controlled conditions for the first time.

\paragraph{Parameter-Efficient Fine-Tuning for Code.}
LoRA~\cite{hu2022lora} enables low-rank adaptation of large models with minimal trainable parameters; QLoRA~\cite{dettmers2024qlora} extends this to quantized models.
Both have become standard for code-specific fine-tuning, with DeepSeek-Coder~\cite{guo2024deepseekcoder} demonstrating that adapter-based training preserves pretrained code understanding while specializing for downstream tasks.
We use LoRA (rank 16) for all experiments to ensure that capacity differences between formulations reflect the task structure rather than total trainable parameters. Our rank ablation (Appendix~\ref{app:ablations}) confirms that doubling adapter capacity does not resolve the generation bottleneck, consistent with a formulation-level constraint rather than a capacity limitation.

\paragraph{Classification vs.\ Generation in NLP.}
Task framing affects model performance broadly: T5~\cite{raffel2020t5} and BART~\cite{lewis2020bart} demonstrate generation-classification tradeoffs under controlled capacity, while SQuAD~\cite{rajpurkar2016squad} shows that constraining output space can simplify tasks without sacrificing quality. We extend this lens to code verification, finding that code complexity moderates the tradeoff in a model-contingent manner.
